\scanpage{34}
sont les sections de \unsure{ramification} de $P^1_S$ sur $P^1_S/(\mathbb Z/3)$. Passant au revêtement galoisien $S'=S(\sqrt{-3})$ de $S$, d'où \unsure{dans|ds} $P^1_{S'}$ \deleted{\unsure{un domaine de}} \inserted{deux sections} \deleted{\unsure{descente}} par $\overline{S'}$, soient $x'$ et $x''$, qui donnent sur $P^1_{S'}/(\mathbb Z/3)$ deux sections; prenons pour l'une \deleted{la} \inserted{comme} section unité ($\infty$ et $0$ étant déjà choisis ci-dessus) --- tandis que l'autre est $-1$. Cela définit un isomorphisme $(P^1_S/\mathbb Z/3)\times_S S'\simeq P^1_{S'}$.

Ainsi, $P^1_S/(\mathbb Z/3)$ est can. isomorphe à un fibré principal de groupe $G_{mS}$, [associé à un faisceau inversible $\Lambda_S$], $\Lambda_S$ étant muni d'une bi-section \inserted{non ramifiée} symétrique sans ramification, ou (ce qui revient au même) d'une forme quadratique non dégénérée; la \unsure{fibration} de cette bisection n'étant autre que celle du revêtement $S'=S(\sqrt{-3})$.

Passant à $S'=S(\sqrt{-3})$, on obtient le diagramme suivant
\mspage{33}
